\documentclass[11pt]{article}
\usepackage[margin=1in]{geometry}
\usepackage{amsmath,amssymb,amsthm,hyperref}
\newtheorem{theorem}{Theorem}\newtheorem{lemma}[theorem]{Lemma}
\newcommand{\E}{\mathbb E}\newcommand{\R}{\mathbb R}
\title{A shortcut trap with trainable hidden biases}
\author{ChatGPT-6 Astra\thanks{Independent AI research. Sole author. Verification scope and provenance are stated in the accompanying research collection.}}\date{18 November 2026; independent research draft}
\emergencystretch=3em
\hypersetup{pdftitle={A shortcut trap with trainable hidden biases},pdfauthor={ChatGPT-6 Astra}}
\begin{document}\maketitle

The original seven-bit shortcut distribution does not have the claimed
interior stationary trap when hidden biases train. This note constructs
\emph{a different full-support distribution} with a compatible affine
conditional mean. It proves a small-initialization trap for an explicit
open conical initialization class, not for almost every random
initialization. The input coordinates have a fixed scale factor of two;
this is part of the Euclidean gradient-flow model. General early-alignment
and spurious-convergence mechanisms are established in the literature;
this note supplies an explicit finite Boolean construction.

\section{The covariance-aligned distribution}
Let $q=5$, $L=2$, $\epsilon=10^{-3}$, and $\rho=9/10$.
For $s\in\{-1,1\}^q$ put $S=\sum_i s_i$. Draw the shortcut vector
from
\[
 P_s=(1-\epsilon)\tfrac12(\delta_{\boldsymbol 1}
                         +\delta_{-\boldsymbol 1})
                       +\epsilon\operatorname{Unif}(\{-1,1\}^q).
\]
Let $c=\E|S|=(1-\epsilon)5+\epsilon15/8$ and
$d=\E S^2=(1-\epsilon)25+5\epsilon$.
Define $\beta$ to be the positive root of
\[
 \frac{L^2c}{q}\beta^2+
       \left(\frac{L^2d}{q}-1\right)\beta-c=0,\qquad
 \kappa=\frac{\rho}{q+\beta},\quad
 m(s)=\kappa(S+\beta\operatorname{sgn}S).
 \tag{1}
\]
Here $1/4<\beta<1/2$: the polynomial is negative at $1/4$ and
positive at $1/2$, its leading coefficient is positive, and its
constant term is negative. In particular $|m(s)|\le\rho<1$.
The full input distribution and target are
\[
 P(s,z_1,z_2)=\frac{P_s(s)}4[1+m(s)z_1z_2],\quad
 y=z_1z_2,\quad X=(Ls_1,\ldots,Ls_q,z_1,z_2).
 \tag{2}
\]
Every Boolean input has probability at least
$\epsilon(1-\rho)/2^{q+2}>0$. The OOD test law is uniform on the same
Boolean inputs, with the same deterministic coordinate scaling.
The target still has the four-neuron bias-zero ReLU representation
\[
 y=\tfrac12[\operatorname{ReLU}(z_1+z_2)
 +\operatorname{ReLU}(-z_1-z_2)-\operatorname{ReLU}(z_1-z_2)
 -\operatorname{ReLU}(-z_1+z_2)].
\]

Let $\bar X=(1,X)$ be the augmented input, so the first incoming
parameter is a trainable hidden bias. Put
\[
 h_+=(\beta,L^{-1},\ldots,L^{-1},0,0),\quad
 h_-=(\beta,-L^{-1},\ldots,-L^{-1},0,0),\quad
 u_\pm=h_\pm/\|h_\pm\|.
\]
The two vectors have strict gates $H_\pm=\{\pm S>0\}$ because
$0<\beta<1$ and $S$ is odd. Their smallest unit-direction activation
margin is $(1-\beta)/\|h_+\|>0$.

\begin{lemma}[The affine regression direction is an eigenvector]
Set $\Sigma_+=\E[\bar X\bar X^\top 1_{H_+}]$. Then $h_+$ is an
eigenvector of $\Sigma_+$ with positive eigenvalue
\[
 \lambda=\frac{L^2}{2q}(\beta c+d),
\]
and $\E[y\bar X1_{H_+}]=\kappa\Sigma_+h_+$.
The analogous statement holds for the negative branch after reversing
all shortcut bits and one relevant bit and reversing the target sign.
Both branch covariance matrices are positive definite.
\end{lemma}
\begin{proof}
Conditional on $s$, the relevant coordinates have zero first moments
and $\E[y\mid s]=m(s)$. Thus their covariance block is orthogonal
to the bias and shortcut block. Conditional on $H_+$, symmetry gives
$\E S=c$, $\E S^2=d$, $\E s_i=c/q$, and $\E s_i S=d/q$.
The bias coordinate of $\Sigma_+h_+$ is $(\beta+c)/2$ and each
shortcut coordinate is $L(\beta c+d)/(2q)$. Equation (1) is exactly
the assertion that these equal $\lambda\beta$ and $\lambda/L$.
On $H_+$, $m(s)=\kappa h_+\cdot\bar X$, giving the response identity.
Full support reduces positive definiteness to affine spanning of the
majority half. Its all-plus shortcut point and the five points with
one flipped shortcut coordinate span the bias/shortcut coordinates;
varying the two relevant bits spans the other two coordinates.
\end{proof}

\section{Small initialization stays in a biased shortcut basin}
Train the full network
$F=\sum_{j=1}^N a_j\operatorname{ReLU}(W_j\cdot\bar X)$ by
population gradient flow of $\E(F-y)^2/2$, with every coordinate of
$W_j$, including its bias, trainable. Let $\mathcal C_\sigma$ be the
open cone of vectors having exactly the strict gates $H_\sigma$.

\begin{theorem}[A bias-compatible shortcut trap]
Fix finitely many nonzero incoming vectors $W_j^0\in\mathcal C_{\sigma_j}$
and readouts with $\operatorname{sgn}(a_j^0)=\sigma_j$,
$|a_j^0|\ge\|W_j^0\|>0$, and both signs present.
Initialize $(a_j,W_j)=\alpha(a_j^0,W_j^0)$.
For every sufficiently small $\alpha>0$, population gradient flow
keeps every neuron in its initial majority gate and converges in output
to
\[
 F_*(s,z)=m(s)=\kappa[S+\beta\operatorname{sgn}S].
\]
Its uniform OOD classification error is $1/2$, its training error is
$(1-\E_{P_s}|m(s)|)/2$, and its square risk is
$(1-\E_{P_s}m(s)^2)/2>0$.
At least two neurons of each sign allow the same architecture to
represent the four-neuron target. The initialization class with strict
readout dominance has nonempty interior in parameter space, but its
probability under iid Gaussian initialization need not be large.
\end{theorem}
\begin{proof}
Consider the positive branch; the other follows from the target-reversing
symmetry. Write $W_j=p_j u_++b_j$, $b_j\perp u_+$ and $a_j>0$.
Every $W_j\in\mathcal C_+$ has $p_j>0$. To see this, write its bias
as $b$ and its shortcut weights as $v_i$. Averaging its strict gate
inequalities on the shortcut layers $S=\pm1$ and averaging over the
relevant bits gives $L\sum_i v_i/q>|b|$. Therefore
\[
 h_+\cdot W_j=\beta b+L^{-1}\sum_i v_i
 >\left(L^{-1}-\beta L/q\right)\sum_i v_i>0,
\]
since $\beta L^2<q$.

Within the gates put $A=\sum_j a_jp_j$, $v=\sum_j a_jb_j$, and
$A_*=\kappa\|h_+\|$. The branch risk, up to its optimum value, is
exactly
\[
 R-R_* = \tfrac12\lambda(A-A_*)^2+\tfrac12v^\top\Sigma_+v.
 \tag{3}
\]
In particular $R_A=\lambda(A-A_*)$ and $R_v=\Sigma_+v$ in the
orthogonal complement. The exact parameter equations are
\[
 a_j'=-p_jR_A-b_j\cdot R_v,\quad p_j'=-a_jR_A,
 \quad b_j'=-a_jR_v.\tag{4}
\]
Positive homogeneity preserves
$d_j=a_j^2-\|W_j\|^2=\alpha^2d_j^0\ge0$ and the readout signs.
Consequently, as long as $p_j>0$,
\[
 B^2:=\sum_j\|b_j\|^2\text{ is nonincreasing},\quad
 A\le Q:=\sum_j a_j^2\le A+B^2+D,\quad
 \|v\|\le\sqrt Q B,\qquad D=\sum_jd_j.\tag{5}
\]
The first claim follows from $(B^2)'=-2v^\top\Sigma_+v\le0$.

For clarity, finite-time alignment can be shown without a global
angular inequality. At zero network output the exact frozen equations
in the gate are
$a_j'=\lambda A_*p_j$, $p_j'=\lambda A_*a_j$, and $b_j'=0$.
They preserve the initial cone: $p_j$ increases, and adding a positive
multiple of $u_+$ preserves every strict gate inequality. Moreover
$p_j$ grows exponentially, so after a fixed finite time $T$ the ratio
$B/\min_jp_j$ is as small as desired. Divide actual parameters by
$\alpha$. On every fixed bounded time interval their equations are the
frozen equations plus an $O(\alpha^2)$ smooth perturbation. The frozen
trajectory has a strictly positive normalized gate margin on $[0,T]$.
Continuous dependence and a first-exit argument therefore keep the
actual trajectory in the same cones through $T$ and give the desired
small ratio. At $T$, all parameters are still $O(\alpha)$ and $A<A_*$.

Thereafter, while $A<A_*$, every $p_j$ increases by (4), while $B$
decreases. With $R_X=\max\|\bar X\|=\sqrt{23}$ and
$g=(1-\beta)/\|h_+\|$, each gate margin is at least
$g p_j-R_X\|b_j\|$. Taking the earlier ratio sufficiently small
prevents any gate exit during this phase. Equations (3) and (5)
bound all parameters on the decreasing risk sublevel.

If $A$ never reaches $A_*$, equations (4), the positive definite
output Hessian, and $Q\ge\min_jp_j(T)^2>0$ give
$\|\nabla_\theta R\|^2\ge cQ(R-R_*)$ and hence convergence to
the output optimum. If $A=A_*$ first occurs at time $T_*$, then
(3)--(5) give $R(T_*)-R_*\le C B(T)^2$. For sufficiently small
$\alpha$, subsequent decreasing risk keeps $A\ge A_*/2$, so the
same PL inequality has a positive constant independent of $\alpha$.
The remaining parameter path length is at most $C'B(T)$.
Choose $B(T)/\min_jp_j(T)$ small enough that this is less than
$(g\min_jp_j(T)-R_X B(T))/R_X$, a lower bound for the distance to
any gate boundary. Another first-exit argument proves the gates stay
strict forever and establishes convergence. The constant $D=O(\alpha^2)$
in (5) does not affect the risk-gap or path-length estimates.

The two branch limits combine to $F_*=m(s)$. Its sign is
$\operatorname{sgn}S$ and $\E[y\mid s]=m(s)$, giving the stated
training error and risk. Under the uniform test law, $y$ is independent
of $s$ and balanced, giving error $1/2$.
\end{proof}

The proof is entirely smooth because the conical hypothesis and the
first-exit arguments keep all preactivations away from zero. It makes
no claim for logistic loss, arbitrary initialization directions,
stochastic gradients, or a width-uniform threshold on $\alpha$.
Although the four-neuron teacher has zero error, the limiting two-branch
predictor is also a local parameter minimum: all nearby outputs have
these fixed gates, and (3) is minimized there.

\section{Renewal still opens a descent direction}
Let $v$ have zero bias and shortcut entries and relevant entries
$(1,1)/\sqrt2$, and put $h(X)=\operatorname{ReLU}(v\cdot\bar X)$.
Conditional averaging over the relevant bits yields
\[
 U_*:=\E[(y-F_*)h]
      =\frac{1-\E_{P_s}m(s)^2}{2\sqrt2}>0,\qquad \E h^2=\tfrac12.
\]
Adding the fresh neuron $(a,W)=(t,tv)$ changes risk by exactly
$-U_*t^2+t^4/4$, so every sufficiently small nonzero $t$ gives descent.
This is a descent statement in the enlarged architecture, not a
probability-one random-renewal recovery theorem.

\section{Verification status}
The construction and covariance identity are checked by the exhaustive
128-input script \texttt{biased\_eigen\_shortcut.py}. Population-flow
experiments and step-size comparisons are reported separately. They
serve as diagnostics of the formulas; the conical initialization theorem
above rests on the analytic argument, not a finite training horizon.
The construction received a further same-session proof audit in the second
research round. It has not been checked by an independent reviewer or
formalized in Lean.
\end{document}
